\documentclass[preview]{standalone}

\usepackage{amsmath}
\usepackage{amssymb}
\usepackage{stellar}
\usepackage{definitions}

\begin{document}

\id{commutative-ring-ideals}
\genpage

\section{Principal, prime, and maximal ideals}

\begin{snippetdefinition}{principal-ideal-definition}{Principal ideals and principal ideal domains}
    Let \(A\) be a \commutativering and let \(x\in A\). The \emph{principal
    ideal generated by \(x\)} is
    \[
        (x)=\{ax\suchthat a\in A\}.
    \]
    An \integraldomain is a \emph{principal ideal domain}, abbreviated
    \emph{PID}, if each of its \ideal[ideals] is principal.
\end{snippetdefinition}

\begin{snippetdefinition}{prime-ideal-definition}{Prime ideal}
    Let \(A\) be a \commutativering. A proper \ideal \(I\idealin A\) is
    \emph{prime} if, for all \(a,b\in A\),
    \[
        ab\in I
        \implies
        a\in I\ \lor\ b\in I.
    \]
\end{snippetdefinition}

\begin{snippetdefinition}{maximal-ideal-definition}{Maximal ideal}
    Let \(A\) be a \commutativering. A proper \ideal \(I\idealin A\) is
    \emph{maximal} if every \ideal \(J\idealin A\) satisfying
    \[
        I\subseteq J\subseteq A
    \]
    is equal to \(I\) or to \(A\).
\end{snippetdefinition}

\begin{snippetexample}{integer-prime-maximal-ideals-example}{Prime and maximal ideals of the integers}
    If \(p\) is a \primen, then
    \[
        (p)=p\integers
    \]
    is both a \primeideal and a \maximalideal of \(\integers\).
    Indeed, Euclid's lemma shows that
    \(p\divides ab\) implies \(p\divides a\) or \(p\divides b\).
    Moreover, if
    \[
        (p)\subseteq(n)\subseteq\integers,
    \]
    then \(n\divides p\), so \((n)=(p)\) or \((n)=\integers\).
    Distinct positive primes therefore determine distinct maximal ideals.
\end{snippetexample}

\begin{snippetproposition}{principale-ideal-domain-every-prime-ideal-is-maximal}{Nonzero prime ideals in a PID are maximal}
    Let \(A\) be a \principalidealdomain. Every nonzero \primeideal of
    \(A\) is a \maximalideal.
\end{snippetproposition}

\begin{snippetproof}{principale-ideal-domain-every-prime-ideal-is-maximal-proof}{principale-ideal-domain-every-prime-ideal-is-maximal}{Nonzero prime ideals in a PID are maximal}
    Let \(I=(x)\) be a nonzero \primeideal, and suppose
    \[
        I\subseteq J\subseteq A.
    \]
    Since \(A\) is a PID, \(J=(y)\) for some \(y\in A\). The inclusion
    \(x\in(y)\) gives \(x=ay\) for some \(a\in A\). Since \(ay=x\in I\)
    and \(I\) is prime, either \(y\in I\) or \(a\in I\).

    If \(y\in I\), then \(J=(y)\subseteq I\), so \(J=I\). Otherwise
    \(a=cx\) for some \(c\in A\), and
    \[
        x=ay=cxy.
    \]
    Hence \(x(1_A-cy)=0_A\). Since \(A\) is an \integraldomain and
    \(x\neq0_A\), one has \(cy=1_A\). Thus \(y\) is a unit and \(J=A\).
\end{snippetproof}

\section{Operations on ideals}

\begin{snippetdefinition}{ideal-addition-multiplication-definition}{Ideal addition and multiplication}
    Let \(A\) be a \commutativering and let \(I,J\idealin A\). Their
    \emph{sum} and \emph{product} are
    \[
        I+J=\{x+y\suchthat x\in I,\ y\in J\}
    \]
    and
    \[
        IJ
        =
        \left\{
            \sum_{h=1}^{n}x_hy_h
            \suchthat
            x_h\in I,\ y_h\in J,\ n\in\naturalnumbers
        \right\}.
    \]
\end{snippetdefinition}

\begin{snippet}{ideal-product-finite-sums-explanation}
    The product \(IJ\) consists of all finite sums of products
    \(xy\), with \(x\in I\) and \(y\in J\); in general it is not merely
    the \set of the individual products \(xy\).
\end{snippet}

\begin{snippetproposition}{ideal-addition-multiplication-is-ideal}{Sums and products of ideals are ideals}
    Let \(A\) be a \commutativering and let \(I,J\idealin A\). Then
    \[
        I+J\idealin A,
        \qquad
        IJ\idealin A.
    \]
\end{snippetproposition}

\begin{snippetproof}{ideal-addition-multiplication-is-ideal-proof}{ideal-addition-multiplication-is-ideal}{Sums and products of ideals are ideals}
    The \set \(I+J\) contains \(0_A\). If
    \(z_k=x_k+y_k\), with \(x_k\in I\) and \(y_k\in J\), then
    \[
        z_1-z_2=(x_1-x_2)+(y_1-y_2)\in I+J.
    \]
    For \(a\in A\),
    \[
        az_1=ax_1+ay_1\in I+J.
    \]
    Thus \(I+J\) is an \ideal.

    The zero element belongs to \(IJ\). If
    \[
        z_1=\sum_{h=1}^{r}x_hy_h,
        \qquad
        z_2=\sum_{k=1}^{s}v_kw_k,
    \]
    then
    \[
        z_1-z_2
        =
        \sum_{h=1}^{r}x_hy_h
        +
        \sum_{k=1}^{s}(-v_k)w_k
        \in IJ.
    \]
    Moreover,
    \[
        az_1=\sum_{h=1}^{r}(ax_h)y_h\in IJ.
    \]
    Therefore \(IJ\) is also an \ideal.
\end{snippetproof}

\begin{snippetdefinition}{ideal-generator-definition}{Finitely generated ideal}
    Let \(A\) be a \commutativering and let
    \(x_1,\ldots,x_n\in A\). The \emph{ideal generated by}
    \(x_1,\ldots,x_n\) is
    \[
        (x_1,\ldots,x_n)
        =
        \left\{\sum_{i=1}^{n}a_ix_i
        \suchthat a_1,\ldots,a_n\in A\right\}
        =
        (x_1)+\cdots+(x_n).
    \]
    An \ideal is \emph{finitely generated} if it has this form.
\end{snippetdefinition}

\begin{snippetproposition}{finitely-generated-ideal-smallest-containing}{The ideal generated by finitely many elements}
    Let \(A\) be a \commutativering and let \(x_1,\ldots,x_n\in A\).
    The \ideal \((x_1,\ldots,x_n)\) is the smallest \ideal of \(A\)
    containing every \(x_i\).
\end{snippetproposition}

\begin{snippetproof}{finitely-generated-ideal-smallest-containing-proof}{finitely-generated-ideal-smallest-containing}{The ideal generated by finitely many elements}
    The \set \((x_1,\ldots,x_n)\) is an \ideal and contains each \(x_i\).
    If an \ideal \(J\) contains all the \(x_i\), then it contains every
    sum \(\sum_i a_ix_i\). Hence
    \((x_1,\ldots,x_n)\subseteq J\).
\end{snippetproof}

\begin{snippetexample}{integer-ideal-product-example}{A product of ideals in the integers}
    In \(\integers\),
    \[
        (3)(5)=(15).
    \]
    Every finite sum of products \((3u_h)(5v_h)\) is a multiple of \(15\),
    so \((3)(5)\subseteq(15)\). Conversely,
    \(15k=(3k)5\in(3)(5)\), proving the reverse inclusion. For example,
    \(30=(6)5\in(3)(5)\).
\end{snippetexample}

\section{Ideals and quotient rings}

\begin{snippetproposition}{prime-maximal-ideals-quotient-characterization}{Prime and maximal ideals through quotients}
    Let \(A\) be a \commutativering and let \(I\idealin A\) be proper.
    Then:
    \begin{enumerate}
        \item \(I\) is a \primeideal \ifandonlyif \(A/I\) is an \integraldomain;
        \item \(I\) is a \maximalideal \ifandonlyif \(A/I\) is a \field.
    \end{enumerate}
\end{snippetproposition}

\begin{snippetproof}{prime-maximal-ideals-quotient-characterization-proof}{prime-maximal-ideals-quotient-characterization}{Prime and maximal ideals through quotients}
    If \(I\) is prime and
    \((a+I)(b+I)=0_{A/I}\), then \(ab\in I\), so \(a\in I\) or
    \(b\in I\). Hence \(A/I\) has no nonzero zero divisors. Conversely,
    if \(A/I\) is an \integraldomain and \(ab\in I\), then
    \[
        (a+I)(b+I)=0_{A/I},
    \]
    whence \(a\in I\) or \(b\in I\). This proves the first equivalence.

    Suppose \(I\) is maximal and \(a+I\neq0_{A/I}\). Then \(a\notin I\),
    so the \ideal \(I+(a)\) strictly contains \(I\) and must equal \(A\).
    Thus \(1_A=x+ab\) for some \(x\in I\) and \(b\in A\), giving
    \[
        (a+I)(b+I)=1_A+I.
    \]
    Therefore \(A/I\) is a \field.

    Conversely, suppose \(A/I\) is a \field and
    \(I\subsetneq J\idealin A\). Choose \(a\in J\difference I\). The class
    \(a+I\) has an inverse \(b+I\), so \(ab-1_A\in I\subseteq J\).
    Since \(ab\in J\), it follows that \(1_A\in J\), and hence \(J=A\).
\end{snippetproof}

\begin{snippetexample}{quotient-ring-new-units-example}{New units in a quotient ring}
    In \(\integers/5\integers\),
    \[
        [2]_5[3]_5=[1]_5.
    \]
    Thus \([2]_5\) is invertible even though \(2\) is not a unit of
    \(\integers\). Invertibility can therefore appear after passing to a
    quotient.
\end{snippetexample}

\begin{snippettheorem}{maximal-ideal-existence}{Existence of maximal ideals}
    Every nontrivial \commutativering has a \maximalideal.
\end{snippettheorem}

\begin{snippetproof}{maximal-ideal-existence-proof}{maximal-ideal-existence}{Existence of maximal ideals}
    Let \(A\) be a nontrivial \commutativering, and order the nonempty
    \set
    \[
        \mathcal{S}=\{I\idealin A\suchthat I\neq A\}
    \]
    by inclusion. If \(\mathcal{C}\subseteq\mathcal{S}\) is a chain, then
    \[
        J=\bigcup_{I\in\mathcal{C}}I
    \]
    is an \ideal: any two of its elements lie together in one member of
    the chain, which supplies closure under differences, and absorption is
    immediate. Moreover \(J\neq A\), since \(1_A\in J\) would place
    \(1_A\) in some proper \ideal in \(\mathcal{C}\). Thus every chain has
    an upper bound in \(\mathcal{S}\). Zorn's lemma gives a maximal element
    of \(\mathcal{S}\), which is a \maximalideal of \(A\).
\end{snippetproof}

\end{document}
